% GENERATED FILE. Edit canonical lessons, then run npm run generate:books.
% canonical-source-hash: fnv1a64:4516f81419110990
% canonical-lessons: UR-C117-badam, UR-C117-mungphali, UR-C117-dahi, UR-C117-tel, UR-C117-ata

\chapter{Almond, Peanut, Yoghurt, Oil, Flour}
\label{ch:ur-almond-peanut-yoghurt-oil-flour}
\hlchaptermodality{117}

\begin{chapteropening}
\textbf{By the end of this chapter:} \emph{I can say an almond, a peanut, yoghurt, oil and flour.}

Complete the last lesson of chapter 117: I can say an almond, a peanut, yoghurt, oil and flour.
\end{chapteropening}

\section[\texorpdfstring{\ur{بادام} (bādām)}{بادام (bādām)}]{\ur{بادام} (bādām) — an almond}

\label{lesson:UR-C117-badam}



\begin{quote}
Before the new one: say the Urdu for a pear, then the Urdu for a date.
\end{quote}

\subsection*{You'll want to know: \ur{بادام}}
\textbf{\ur{بادام}} — \emph{bādām} — \textquotedblleft{}an almond\textquotedblright{}.

\begin{cousinweb}
\textbf{\ur{بادام}} (\emph{bādām}) is a loan, usually traced to Persian \emph{bādām}, \textquotedblleft{}almond\textquotedblright{}.
\end{cousinweb}

\begin{grammarlens}[title={What you've built}]
One more word for this chapter.
\end{grammarlens}

\subsection*{Your turn}
\begin{itemize}
\raggedright
  \item \emph{Say it:} \emph{bādām}
  \item \emph{Say it:} \emph{bādām}, once more
  \item \emph{Say it:} say \emph{bādām}
  \item \emph{Recall:} say the Urdu for a pear, then the Urdu for a date, then say \emph{bādām} again
\end{itemize}

\begin{tcolorbox}[breakable,colback=teal!4,colframe=teal!35!black,title={Before you move on}]
What does \ur{بادام} mean? (\textquotedblleft{}An almond\textquotedblright{}.) Say it once more.
\end{tcolorbox}
\section[\texorpdfstring{\ur{مونگ} \ur{پھلی} (mūṅgphalī)}{مونگ پھلی (mūṅgphalī)}]{\ur{مونگ} \ur{پھلی} (mūṅgphalī) — a peanut}

\label{lesson:UR-C117-mungphali}



\begin{quote}
Before the new one: say the Urdu for a date, then the Urdu for an almond.
\end{quote}

\subsection*{You'll want to know: \ur{مونگ} \ur{پھلی}}
\textbf{\ur{مونگ} \ur{پھلی}} — \emph{mūṅgphalī} — \textquotedblleft{}a peanut\textquotedblright{}.

\begin{cousinweb}
\textbf{\ur{مونگ} \ur{پھلی}} (\emph{mūṅgphalī}) is \textquotedblleft{}mung-pod\textquotedblright{}: \emph{mūṅg}, \textquotedblleft{}the mung bean\textquotedblright{}, from Sanskrit \emph{mudga} (CDIAL 10198), and \emph{phalī}, \textquotedblleft{}pod\textquotedblright{}, from \emph{phala}, \textquotedblleft{}fruit\textquotedblright{} (CDIAL 9051).
\end{cousinweb}

\begin{grammarlens}[title={What you've built}]
One more word for this chapter.
\end{grammarlens}

\subsection*{Your turn}
\begin{itemize}
\raggedright
  \item \emph{Say it:} \emph{mūṅgphalī}
  \item \emph{Say it:} \emph{mūṅgphalī}, once more
  \item \emph{Say it:} say \emph{mūṅgphalī}
  \item \emph{Recall:} say the Urdu for a date, then the Urdu for an almond, then say \emph{mūṅgphalī} again
\end{itemize}

\begin{tcolorbox}[breakable,colback=teal!4,colframe=teal!35!black,title={Before you move on}]
What does \ur{مونگ} \ur{پھلی} mean? (\textquotedblleft{}A peanut\textquotedblright{}.) Say it once more.
\end{tcolorbox}
\section[\texorpdfstring{\ur{دہی} (dahī)}{دہی (dahī)}]{\ur{دہی} (dahī) — yoghurt}

\label{lesson:UR-C117-dahi}



\begin{quote}
Before the new one: say the Urdu for an almond, then the Urdu for a peanut.
\end{quote}

\subsection*{You'll want to know: \ur{دہی}}
\textbf{\ur{دہی}} — \emph{dahī} — \textquotedblleft{}yoghurt\textquotedblright{}.

\begin{cousinweb}
\textbf{\ur{دہی}} (\emph{dahī}) is a tadbhava, a word inherited from Sanskrit and reshaped by sound change. It comes from \emph{dadhi}, \textquotedblleft{}thick sour milk\textquotedblright{} (CDIAL 6146).
\end{cousinweb}

\begin{grammarlens}[title={What you've built}]
One more word for this chapter.
\end{grammarlens}

\subsection*{Your turn}
\begin{itemize}
\raggedright
  \item \emph{Say it:} \emph{dahī}
  \item \emph{Say it:} \emph{dahī}, once more
  \item \emph{Say it:} say \emph{dahī}
  \item \emph{Recall:} say the Urdu for an almond, then the Urdu for a peanut, then say \emph{dahī} again
\end{itemize}

\begin{tcolorbox}[breakable,colback=teal!4,colframe=teal!35!black,title={Before you move on}]
What does \ur{دہی} mean? (\textquotedblleft{}Yoghurt\textquotedblright{}.) Say it once more.
\end{tcolorbox}
\section[\texorpdfstring{\ur{تیل} (tel)}{تیل (tel)}]{\ur{تیل} (tel) — oil}

\label{lesson:UR-C117-tel}



\begin{quote}
Before the new one: say the Urdu for a peanut, then the Urdu for yoghurt.
\end{quote}

\subsection*{You'll want to know: \ur{تیل}}
\textbf{\ur{تیل}} — \emph{tel} — \textquotedblleft{}oil\textquotedblright{}.

\begin{cousinweb}
\textbf{\ur{تیل}} (\emph{tel}) is a tadbhava, a word inherited from Sanskrit and reshaped by sound change. It comes from \emph{taila}, \textquotedblleft{}sesame oil, oil\textquotedblright{} (CDIAL 5958). It is built on \emph{tila}, \textquotedblleft{}sesame\textquotedblright{}.
\end{cousinweb}

\begin{grammarlens}[title={What you've built}]
One more word for this chapter.
\end{grammarlens}

\subsection*{Your turn}
\begin{itemize}
\raggedright
  \item \emph{Say it:} \emph{tel}
  \item \emph{Say it:} \emph{tel}, once more
  \item \emph{Say it:} say \emph{tel}
  \item \emph{Recall:} say the Urdu for a peanut, then the Urdu for yoghurt, then say \emph{tel} again
\end{itemize}

\begin{tcolorbox}[breakable,colback=teal!4,colframe=teal!35!black,title={Before you move on}]
What does \ur{تیل} mean? (\textquotedblleft{}Oil\textquotedblright{}.) Say it once more.
\end{tcolorbox}
\section[\texorpdfstring{\ur{آٹا} (āṭā)}{آٹا (āṭā)}]{\ur{آٹا} (āṭā) — flour}

\label{lesson:UR-C117-ata}



\begin{quote}
Before the new one: say the Urdu for yoghurt, then the Urdu for oil.
\end{quote}

\subsection*{You'll want to know: \ur{آٹا}}
\textbf{\ur{آٹا}} — \emph{āṭā} — \textquotedblleft{}flour\textquotedblright{}.

\begin{cousinweb}
\textbf{\ur{آٹا}} (\emph{āṭā}) is a tadbhava, a word inherited and reshaped by sound change. It comes from \emph{*ārta}, \textquotedblleft{}flour\textquotedblright{}, a form CDIAL reconstructs from its descendants (CDIAL 1338). CDIAL compares Persian \emph{ārd}, \textquotedblleft{}flour\textquotedblright{}.
\end{cousinweb}

\begin{grammarlens}[title={What you've built}]
One more word for this chapter.
\end{grammarlens}

\subsection*{Your turn}
\begin{itemize}
\raggedright
  \item \emph{Say it:} \emph{āṭā}
  \item \emph{Say it:} \emph{āṭā}, once more
  \item \emph{Say it:} say \emph{āṭā}
  \item \emph{Recall:} say the Urdu for yoghurt, then the Urdu for oil, then say \emph{āṭā} again
\end{itemize}

\begin{tcolorbox}[breakable,colback=teal!4,colframe=teal!35!black,title={Before you move on}]
What does \ur{آٹا} mean? (\textquotedblleft{}Flour\textquotedblright{}.) Say it once more.
\end{tcolorbox}
